\documentclass[10pt,a4paper]{article}
\usepackage[margin=0.5in]{geometry}
\usepackage{enumitem}
\usepackage{hyperref}
\usepackage{titlesec}
\usepackage{xcolor}
\usepackage{fontawesome5}
\setlength{\parindent}{0pt}

% Color definitions
\definecolor{darkblue}{RGB}{0,51,102}
\definecolor{mediumgray}{RGB}{80,80,80}

% Hyperref setup
\hypersetup{
    colorlinks=true,
    linkcolor=darkblue,
    urlcolor=darkblue,
    pdftitle={CV - Musthofa Joko Anggoro},
    pdfauthor={Musthofa Joko Anggoro}
}

% Section formatting
\titleformat{\section}{\large\bfseries\color{darkblue}}{}{0em}{}[\titlerule]
\titlespacing*{\section}{0pt}{6pt}{3pt}

\titleformat{\subsection}{\normalsize\bfseries}{}{0em}{}
\titlespacing*{\subsection}{0pt}{4pt}{1pt}

% Remove page numbering
\pagestyle{empty}

% Custom commands
\newcommand{\contactline}[1]{\small #1}
\newcommand{\experienceheader}[4]{%
    \textbf{#1} \hfill #2 \\
    \textit{#3} \hfill \textit{#4}%
}

\begin{document}

% ============================================================
% HEADER
% ============================================================
\begin{center}
    {\LARGE \textbf{MUSTHOFA JOKO ANGGORO}} \\[2pt]
    \contactline{
        \faEnvelope\ \href{mailto:musthofaja.topa@gmail.com}{musthofaja.topa@gmail.com} \ $|$ \
        \faLinkedin\ \href{https://www.linkedin.com/in/musthofa-joko-anggoro/}{linkedin.com/in/musthofa-joko-anggoro} \ $|$ \
        \faGithub\ \href{https://github.com/topahilangharapan}{github.com/topahilangharapan} \ $|$ \
        \faGlobe\ \href{https://portfolio-topahilangharapan-personal.vercel.app/}{portfolio-topahilangharapan-personal.vercel.app}
    }
\end{center}

\vspace{-8pt}

% ============================================================
% RESEARCH INTERESTS
% ============================================================
\section*{RESEARCH INTERESTS}

My research centers on \textbf{reconfigurable computing and FPGA-based hardware acceleration}, with particular interest in mapping compute-intensive workloads onto resource-constrained programmable devices. Through my undergraduate thesis --- a full RTL implementation of the Mano register-transfer model on a Xilinx Spartan-3AN FPGA in VHDL-93 --- I developed hands-on expertise in datapath design, RTL synthesis, place-and-route closure, and hardware verification. I am drawn to research at the intersection of \textbf{FPGA architectures, edge AI hardware, and reconfigurable systems}, asking how hardware-aware design methodologies and accelerator architectures can enable efficient, low-power inference at the edge. I am enthusiastic about contributing to the Hardware \& Embedded Systems Lab (HESL) under Assoc Prof Lam Siew Kei and advancing the lab's work on FPGA-based accelerators and reconfigurable edge platforms.

\vspace{-4pt}

% ============================================================
% EDUCATION
% ============================================================
\section*{EDUCATION}

\experienceheader{University of Indonesia (UI) --- Faculty of Computer Science (Fasilkom)}{Jakarta, Indonesia}{Bachelor of Computer Science, Information Systems}{Aug 2022 -- Aug 2026 (Expected)}
\begin{itemize}[leftmargin=*, itemsep=1pt, topsep=2pt]
    \item \textbf{GPA: 3.84 / 4.00}
    \item \textbf{Relevant Coursework:} Computer Architecture, Operating Systems, Data Structures \& Algorithms, Introduction to AI \& Data Science, Statistics \& Probability, Data Communication Networks, Database Systems
\end{itemize}

\vspace{2pt}
\experienceheader{Nanyang Technological University (NTU) \& ITS}{Singapore \& Surabaya}{INSPIRASI--NTU Summer Programme 2025 \textnormal{--- \textbf{LPDP Scholarship Awardee}}}{July 2025}
\begin{itemize}[leftmargin=*, itemsep=1pt, topsep=2pt]
    \item Applied software engineering and systems principles to sustainable development challenges; produced a technical report proposing a tiered smart-port automation framework integrating IoT sensor networks and renewable energy utilization
\end{itemize}

\vspace{-4pt}

% ============================================================
% UNDERGRADUATE THESIS
% ============================================================
\section*{UNDERGRADUATE THESIS}

\experienceheader{VHDL-93 Implementation of the Mano Register-Transfer Microoperation Model on FPGA}{}{}{}
\vspace{-8pt}
\begin{itemize}[leftmargin=*, itemsep=1pt, topsep=2pt]
    \item[] \textit{Supervisor:} Prof.\ Petrus Mursanto, University of Indonesia \hfill \textit{Status: Completed --- Paper in preparation for ICACSIS 2026}
\end{itemize}
\begin{itemize}[leftmargin=*, itemsep=1pt, topsep=1pt]
    \item Designed and implemented the Mano register-transfer computer model as a complete, hierarchically structured VHDL-93 (IEEE~1076-1993) datapath on a \textbf{Xilinx Spartan-3AN} FPGA
    \item Conducted a systematic \textbf{gap analysis} mapping 5 undefined textbook constraints to concrete RTL design decisions, bridging theory and physical hardware implementation
    \item Implemented a \textbf{three-tier component hierarchy}: leaf-level primitives (\texttt{reg\_n}, ALU, shifter) $\rightarrow$ functional units $\rightarrow$ full system integration (\texttt{fpga\_top})
    \item Verified all modules via \textbf{ISim functional simulation} testbenches and board-level testing on physical hardware
    \item Achieved \textbf{timing closure at 50~MHz} on the Spartan-3AN device (critical path: 17.05~ns) using XST synthesis engine
    \item Demonstrated live operation via PS/2 keyboard input and \textbf{HD44780 LCD display} output; programmed via JTAG
    \item \textit{Toolchain:} Xilinx ISE~14.x $\cdot$ XST $\cdot$ ISim $\cdot$ JTAG $\cdot$ VHDL-93
\end{itemize}

\vspace{-4pt}

% ============================================================
% PUBLICATIONS
% ============================================================
\section*{PUBLICATIONS}

\begin{itemize}[leftmargin=*, itemsep=2pt, topsep=2pt]
    \item Musthofa Joko Anggoro, Petrus Mursanto. ``VHDL-93 Implementation of the Mano Register-Transfer Microoperation Model on an FPGA.'' \textit{Proceedings of the International Conference on Advanced Computer Science and Information Systems (ICACSIS) 2026}. \textbf{[In preparation]}
\end{itemize}

\vspace{-4pt}

% ============================================================
% TECHNICAL SKILLS
% ============================================================
\section*{TECHNICAL SKILLS}

\textbf{HDL \& Digital Design:} VHDL-93 (primary), RTL synthesis, place-and-route, functional simulation, testbench writing, hierarchical hardware decomposition \\
\textbf{FPGA Tools:} Xilinx ISE~14.x $\cdot$ XST synthesis engine $\cdot$ ISim $\cdot$ JTAG; \textbf{Target Device:} Xilinx Spartan-3AN (XC3S700AN) \\
\textbf{Systems Programming:} Python, Java, C, Go, TypeScript \\
\textbf{Machine Learning:} Scikit-learn, supervised/unsupervised learning, model evaluation \\
\textbf{Infrastructure:} Docker, Kubernetes, CI/CD pipelines, Linux/Unix, Git

\vspace{-4pt}

% ============================================================
% EXPERIENCE
% ============================================================
\section*{EXPERIENCE}

\experienceheader{Software Engineer Intern}{Traveloka (Travel Technology, Singapore HQ)}{Systems Engineering \& Backend Development}{Aug 2025 -- Apr 2026}
\begin{itemize}[leftmargin=*, itemsep=1pt, topsep=2pt]
    \item Engineered a secure Credential Management System with AES-GCM encryption, automated rotation, and audit logging on MongoDB, serving enterprise-scale production traffic
    \item Developed automated flight-reschedule pipeline processing 1,000$+$ daily booking modifications with 99.9\% accuracy, reducing manual intervention by 85\%
    \item Enhanced CI/CD pipelines with Karate-based automated API testing and auto-rollback logic; leveraged GitHub Copilot CLI to accelerate development velocity
\end{itemize}

\vspace{2pt}
\experienceheader{Full Stack Developer Intern}{PT Magna Solusi Indonesia}{Enterprise Application Development}{Jan 2025 -- Jun 2025}
\begin{itemize}[leftmargin=*, itemsep=1pt, topsep=2pt]
    \item Developed full-stack enterprise modules (Spring Boot / Angular), optimized PostgreSQL query performance, and built RESTful microservice APIs handling 10,000$+$ daily requests
\end{itemize}

\vspace{-4pt}

% ============================================================
% LEADERSHIP & ACTIVITIES
% ============================================================
\section*{LEADERSHIP \& ACTIVITIES}

\experienceheader{Vice Person in Charge, UI/UX Division}{COMPFEST --- Indonesia's Largest IT Event (Universitas Indonesia)}{}{Jan 2023 -- Dec 2023}
\begin{itemize}[leftmargin=*, itemsep=1pt, topsep=2pt]
    \item Led a team of 11 designers across 5 design sprints; managed cross-functional collaboration with software engineering and automation divisions to ship the COMPFEST website redesign
\end{itemize}

\vspace{2pt}
\experienceheader{Mentor, Programming Foundations 0}{Faculty of Computer Science, Universitas Indonesia}{}{2024}
\begin{itemize}[leftmargin=*, itemsep=1pt, topsep=2pt]
    \item Mentored first-year students in foundational programming concepts and problem-solving methodologies
\end{itemize}

\end{document}
